\section{Investigación complementaria para defensa: componentes restantes del proyecto}

\subsection{Alcance}

Este informe \textbf{no repite}:

\begin{itemize}
  \item \verb|docs/Personal Research/data-structure-and-canonical-schema-research-report.md|
  \item \verb|docs/Personal Research/qrdqn-research-report.md|
\end{itemize}

Aquí se documentan los componentes que faltaban de cara al tribunal: componentes históricos,
baseline clásico, tuning, scripts auxiliares y laboratorio.

\subsection{1) Mapa de ``qué es cada cosa'' (resto del proyecto)}

\begin{table}[htbp]
\footnotesize
\begin{tabular}{p{3cm}p{3.5cm}p{3cm}p{3.5cm}}
\hline
\textbf{Área} & \textbf{Archivo(s) clave} & \textbf{Qué es} & \textbf{Por qué importa en defensa} \\
\hline
Branch histórico NSL-KDD
  & \verb|src/load_nsl_kdd.py|, \verb|experiments/nslkdd_experiments.md|
  & Línea histórica de benchmark
  & Justifica evolución metodológica del TFG \\
Baseline no-RL
  & \verb|src/baseline_random_forest.py|
  & Comparador supervisado clásico (Random Forest)
  & Permite defender que RL no se eligió ``a ciegas'' \\
Tuning de hiperparámetros
  & \verb|src/tune_hparams.py|
  & Búsqueda con Optuna sobre QRDQN
  & Muestra metodología de optimización reproducible \\
Inferencia Fase~2 (legacy)
  & \verb|scripts/predict_real_traffic.py|
  & Versión antigua del pipeline
  & Sirve para explicar por qué existe \texttt{v2} como ruta mantenida \\
Utilidades de robustez
  & \verb|src/scaling_utils.py|
  & Clipping por percentiles y por z-score
  & Explica mitigación práctica de outliers/domain shift \\
Generación de tráfico de laboratorio
  & \verb|lab/docker/docker-compose.yaml|, \verb|lab/docker/generator/gen_traffic.py|
  & Mini-lab aislado + generador de flujos
  & Demuestra estrategia de experimentación segura \\
Archivo de experimentos
  & \verb|experiments/README.md|, \verb|experiments/cicids2017_qrdqn_experiments.md|
  & Historial mantenido vs histórico
  & Ayuda a separar evidencia actual de contexto histórico \\
\hline
\end{tabular}
\end{table}

\subsection{2) Algoritmos adicionales (aparte de QRDQN)}

\subsubsection{2.1 Random Forest baseline}

\textbf{Dónde está}: \verb|src/baseline_random_forest.py|

\begin{itemize}
  \item Modelo: \texttt{RandomForestClassifier} de scikit-learn.
  \item Configuración actual en código:
    \begin{itemize}
      \item \texttt{n\_estimators=200}
      \item \texttt{max\_depth=None}
      \item \texttt{n\_jobs=-1}
      \item \texttt{random\_state=42}
      \item \texttt{class\_weight=None}
    \end{itemize}
  \item Flujo:
    \begin{enumerate}
      \item Carga dataset (CICIDS2017 o NSL-KDD) ya mapeado al contrato del proyecto.
      \item Entrena RF.
      \item Evalúa con matriz de confusión e informe de clasificación.
      \item Guarda modelo \texttt{.joblib} en \verb|models/|.
    \end{enumerate}
\end{itemize}

\textbf{Mensaje para tribunal}: este baseline es el ``control clásico'' para comparar con RL y
demostrar que la línea RL aporta flexibilidad por función de coste (recompensas), no solo accuracy.

\subsubsection{2.2 Optuna (no es modelo, es algoritmo de búsqueda)}

\textbf{Dónde está}: \verb|src/tune_hparams.py|

\begin{itemize}
  \item Crea un estudio con \texttt{optuna.create\_study(direction="maximize")}.
  \item Objetivo: maximizar \texttt{F1} de ataque.
  \item Espacio de búsqueda actual:
    \begin{itemize}
      \item \texttt{learning\_rate}: \texttt{1e-5} a \texttt{1e-3} (log)
      \item \texttt{batch\_size}: \texttt{\{256, 512, 1024, 2048\}}
      \item \texttt{gradient\_steps}: \texttt{\{10, 50, 100\}}
      \item \texttt{net\_arch}: \texttt{\{[256,128], [512,256], [256,256]\}}
      \item \texttt{gamma}: \texttt{0.95} a \texttt{0.999}
      \item \texttt{train\_freq}: \texttt{\{50,100,200\}}
    \end{itemize}
  \item Resultado se persiste en \verb|runs/optuna/study_<timestamp>.json|.
\end{itemize}

\textbf{Mensaje para tribunal}: existe proceso formal de tuning; no es ``prueba y error manual''.

\subsection{3) Parámetros clave que debes poder defender (resto de scripts)}

\subsubsection{3.1 Validación leave-one-exact-CSV-out}

\textbf{Dónde está}: \verb|src/validate_leave_one_csv_out.py|

Parámetros CLI relevantes:

\begin{itemize}
  \item \texttt{--timesteps} (default 30000)
  \item \texttt{--holdout-csvs} (elige folds concretos)
  \item \texttt{--max-rows-per-csv} (smoke/dev)
  \item \texttt{--predict-batch-size} (default 8192)
\end{itemize}

Qué añade frente a validación simple:

\begin{itemize}
  \item Métricas por fold + agregados (\texttt{mean/std/min/max}).
  \item Métricas globales por suma de confusiones.
  \item Métricas operativas adicionales: \texttt{balanced\_accuracy}, \texttt{specificity},
        \texttt{fpr}, \texttt{fnr}, \texttt{block\_rate}, \texttt{reward\_per\_sample}, tiempos de
        train/eval.
\end{itemize}

\subsubsection{3.2 Check B y Check C (detalles prácticos no obvios)}

\textbf{Dónde está}: \verb|src/validate_checks.py|

\begin{itemize}
  \item \textbf{Check~B (anti-leakage)} define umbral:
        \texttt{leakage\_threshold = baseline\_acc + 0.05}.
  \item \textbf{Check~C (split duro por CSV/día)} reentrena QRDQN con hiperparámetros fijos de
        validación (no exactamente los mismos del train principal).
  \item Usa \texttt{ProgressCallback} para trazabilidad de entrenamiento en validaciones largas.
\end{itemize}

\subsubsection{3.3 Pipeline robusto de inferencia: utilidades de clipping}

\textbf{Dónde está}: \verb|src/scaling_utils.py|

\begin{itemize}
  \item \verb|apply_percentile_clipping(X, p_low, p_high)|: recorta features en bruto antes del
        scaler.
  \item \verb|apply_z_clipping(X_scaled, max_z)|: recorta tras escalar a
        \([-\text{max\_z},\, +\text{max\_z}]\).
\end{itemize}

\textbf{Mensaje para tribunal}: son defensas concretas contra colas extremas y drift de distribución
en tráfico real.

\subsection{4) Funciones clave (resto no cubierto)}

\subsubsection{4.1 \texttt{load\_nsl\_kdd\_binary()}}

\textbf{Archivo}: \verb|src/load_nsl_kdd.py|

Qué hace:

\begin{itemize}
  \item Descarga NSL-KDD con \texttt{kagglehub} si hace falta.
  \item Detecta columnas de etiqueta y dificultad.
  \item Binariza \texttt{normal} vs ataque.
  \item Permite modo canónico (con máscara) o modo legacy one-hot.
  \item Devuelve contrato estándar:
        \texttt{(X\_train, y\_train, X\_test, y\_test, scaler, feature\_names)}.
\end{itemize}

Por qué es defendible: es la pieza que conserva continuidad histórica del proyecto, pero separada
del camino final de Fase~2.

\subsubsection{4.2 \texttt{compute\_truth\_metrics()} en inferencia v2}

\textbf{Archivo}: \verb|scripts/predict_real_traffic_v2.py|

Qué hace:

\begin{itemize}
  \item Detecta columnas de verdad-terreno (\texttt{truth\_y}, \texttt{truth\_label}) cuando
        existen.
  \item Normaliza etiquetas a binario (\texttt{0/1}) y descarta filas no válidas.
  \item Calcula métricas de clasificación
        (\texttt{accuracy}, \texttt{precision/recall/f1}, \texttt{tp/tn/fp/fn}).
  \item Integra estas métricas en el \texttt{metrics.json} final de la ejecución.
\end{itemize}

Por qué importa: te permite defender evaluación controlada en Fase~2 cuando el CSV lleva etiquetas
de referencia.

\subsubsection{4.3 \texttt{gen\_traffic.py} (laboratorio)}

\textbf{Archivo}: \verb|lab/docker/generator/gen_traffic.py|

Qué hace:

\begin{itemize}
  \item Genera dos patrones:
    \begin{enumerate}
      \item múltiples conexiones cortas a puertos cerrados (``scan-like'')
      \item muchas peticiones HTTP con conexión nueva en cada request
    \end{enumerate}
  \item Añade \texttt{jitter} temporal para variabilidad.
\end{itemize}

Por qué importa: justifica cómo produces tráfico reproducible/controlado para pruebas de Fase~2.

\subsection{5) Red neuronal y arquitectura (solo lo complementario)}

Sin repetir teoría QRDQN ya investigada:

\begin{itemize}
  \item El pipeline principal usa \texttt{MlpPolicy} (MLP densa, no CNN/RNN).
  \item El repo también prueba arquitecturas alternativas en tuning
        (\texttt{[256,128]}, \texttt{[256,256]}, \texttt{[512,256]}).
  \item En validaciones y entrenamiento se mantiene acción discreta binaria (PERMIT/BLOCK),
        coherente con salida por acción del modelo.
\end{itemize}

Mensaje defendible: se eligió arquitectura MLP porque el input es vector tabular de flujo, no
imagen/secuencia cruda.

\subsection{6) Dependencias técnicas y rol de cada una}

\textbf{Archivo}: \verb|requirements.txt|

\begin{itemize}
  \item \texttt{numpy}, \texttt{pandas}: procesamiento numérico y tabular.
  \item \texttt{scikit-learn}: scaler, métricas, Random Forest.
  \item \texttt{gymnasium}: API del entorno RL.
  \item \texttt{torch}: backend de redes.
  \item \texttt{stable-baselines3} + \texttt{sb3-contrib}: algoritmos RL (incl. QRDQN).
  \item \texttt{tensorboard}: trazas de entrenamiento.
  \item \texttt{matplotlib}: soporte de visualización.
  \item \texttt{optuna}: tuning automatizado.
\end{itemize}

Mensaje para tribunal: stack estándar y auditable; no hay dependencias exóticas sin justificación.

\subsection{7) Legacy vs mantenido (muy útil en preguntas ``trampa'')}

\begin{itemize}
  \item \textbf{Mantenido para Fase~2}: \verb|scripts/predict_real_traffic_v2.py|
  \item \textbf{Legacy}: \verb|scripts/predict_real_traffic.py| (carga datos de entrenamiento al
        import y no separa artefactos con el mismo nivel de robustez/CLI)
\end{itemize}

Cómo responder si preguntan por ambas: la v1 queda como referencia histórica/técnica; la v2 es la
ruta oficial por diseño reproducible, CLI explícita y endurecimiento frente a domain shift.

\subsection{8) Guion corto para defensa oral (solo ``resto'')}

\begin{enumerate}
  \item ``Además del núcleo canónico + QRDQN, el proyecto tiene un bloque de \textbf{gobernanza
        experimental}: validación dura por folds exactos, agregación de métricas y archivado de
        runs.''
  \item ``Mantengo un baseline clásico (Random Forest) para comparación metodológica, no solo para
        maximizar métrica.''
  \item ``Tengo \textbf{tuning reproducible} con Optuna y búsqueda declarada de hiperparámetros.''
  \item ``La Fase~2 se apoya en un mini-lab aislado y scripts de tráfico controlado, no en pruebas
        opacas.''
  \item ``La inferencia v2 puede incluir métricas con verdad-terreno cuando existen etiquetas,
        mejorando la trazabilidad experimental.''
\end{enumerate}

\subsection{9) Preguntas probables del tribunal y respuesta breve}

\textbf{¿Para qué conservar NSL-KDD si no es el baseline final?}
Para justificar evolución metodológica y comparar decisiones de diseño; el camino mantenido final es
CICIDS2017 + pipeline de Fase~2.

\textbf{¿Por qué incluir Random Forest si el TFG es RL?}
Porque un baseline clásico ancla la comparación y evita conclusiones sin referencia.

\textbf{¿Qué aporta leave-one-exact-CSV-out frente a métricas agregadas?}
Mide robustez por archivo real específico y reduce optimismo por mezcla de patrones entre
train/test.

\textbf{¿Qué pasa si en Fase~2 no hay \texttt{truth\_label} o \texttt{truth\_y}?}
La pipeline sigue funcionando en modo no supervisado operativo y reporta tasas de
\texttt{block/allow}; las métricas supervisadas solo se añaden cuando hay verdad-terreno válida.

\subsection{10) Conclusión de esta investigación complementaria}

Tus dos investigaciones cubren bien el núcleo (contrato de datos + QRDQN). Lo que faltaba para una
defensa sólida era demostrar dominio del \textbf{ecosistema completo}: baseline clásico, tuning,
validación avanzada por fold, tooling de laboratorio y distinción legacy/mantenido. Ese ``resto'' es
exactamente lo que completa el relato técnico ante tribunal.
